\SourcePage{599}{602}
\section{Hyperfine structure of atomic levels}

Another atomic effect associated with specific properties of the nucleus is
the splitting of atomic energy levels through the interaction of the
electrons with the nuclear spin. This is the so-called \emph{hyperfine
structure} of the levels. Since this interaction is weak, the intervals of
the resulting structure are very small, even relative to the fine-structure
intervals. The hyperfine structure must therefore be considered separately
for each fine-structure component.

In this section we denote the nuclear spin by $\mathbf i$ (in accordance with
the convention of atomic spectroscopy), while retaining $\mathbf J$ for the
total angular momentum of the atomic electron shell. Let the total angular
momentum of the atom, including the nucleus, be
$\mathbf F=\mathbf J+\mathbf i$. Each hyperfine-structure component is
characterized by a definite value of this angular momentum. From the general
rules for addition of angular momenta, the quantum number $F$ takes the values
\begin{equation}
 F=J+i,\quad J+i-1,\ldots,|J-i|,
 \tag{121.1}
\end{equation}
so that every level of given $J$ splits into $2i+1$ components if $i<J$, or
into $2J+1$ components if $i>J$.

Because the mean electron distances $r$ in an atom are large compared with
the nuclear radius $R$, the principal contributions to hyperfine splitting
come from the interactions of the electrons with the lowest-order nuclear
multipole moments. These are the magnetic dipole and electric quadrupole
moments (the mean electric dipole moment vanishes; see \S~75).

The nuclear magnetic moment is of order
$\mu_{\text{nuc}}\sim eRv_{\text{nuc}}/c$, where $v_{\text{nuc}}$ denotes
the velocities of nucleons in the nucleus. Its interaction energy with the
electron magnetic moment ($\mu_{\text{el}}\sim e\hbar/mc$) is of order

\SourcePage{600}{603}
\begin{equation}
 \frac{\mu_{\text{nuc}}\mu_{\text{el}}}{r^3}
 \sim\frac{e^2\hbar}{mc^2}\frac{Rv_{\text{nuc}}}{r^3}.
 \tag{121.2}
\end{equation}

The nuclear quadrupole moment is $Q\sim eR^2$; the interaction energy of its
field with the electron charge is of order
\begin{equation}
 \frac{eQ}{r^3}\sim\frac{e^2R^2}{r^3}.
 \tag{121.3}
\end{equation}
Comparison of (121.2) and (121.3) shows that the magnetic interaction, and
hence the level splitting that it produces, exceeds the quadrupole
interaction by a factor
$(v_{\text{nuc}}/c)(\hbar/mcR)\sim15$. Although $v_{\text{nuc}}/c$ is
comparatively small, $\hbar/mcR$ is large.

The operator for the magnetic interaction of the electrons with the nucleus
has the form
\begin{equation}
 \hat V_{iJ}=a\hat{\mathbf i}\hat{\mathbf J}
 \tag{121.4}
\end{equation}
(in analogy with the electron spin--orbit interaction (72.4)). Its dependence
on $F$, and consequently that of the splitting which it produces, is
\begin{equation}
 \frac a2F(F+1)
 \tag{121.5}
\end{equation}
(cf. (72.5)).

The quadrupole interaction operator is formed from the operator
$\hat Q_{ik}$ of the nuclear quadrupole-moment tensor and the components of
the electron angular momentum $\hat{\mathbf J}$. It is proportional to the
scalar $\hat Q_{ik}\hat J_i\hat J_k$ constructed from these operators, and
thus has the form
\begin{equation}
 b\left[\hat i_i\hat i_k+\hat i_k\hat i_i
 -\frac23i(i+1)\delta_{ik}\right]\hat J_i\hat J_k;
 \tag{121.6}
\end{equation}
This result uses the fact that $Q_{ik}$ can be expressed through the nuclear-spin operator by a formula of the form

(75.2). Evaluating the eigenvalues of operator (121.6)---by exactly the same procedure as in Problem~1 of \S~84---gives the following dependence of the quadrupole hyperfine level splitting on the quantum number $F$:

\begin{equation}
 \frac b2F^2(F+1)^2
 +\frac b2F(F+1)[1-2J(J+1)-2i(i+1)].
 \tag{121.7}
\end{equation}

\SourcePage{601}{604}
Magnetic hyperfine splitting is especially pronounced for levels associated
with an outer electron in an $s$ state, because such an electron has a
comparatively high probability of being near the nucleus.

We calculate the hyperfine splitting for an atom having one outer $s$
electron (E.~Fermi, 1930). Its motion in the self-consistent field of the
remaining electrons and the nucleus is described by a spherically symmetric
wavefunction $\psi(r)$\footnote{The calculation below assumes that
$Ze^2/(\hbar c)\ll1$ (cf. the note on p.~601).}.

We seek the electron--nucleus interaction operator as the energy operator
$-\hat{\boldsymbol\mu}\hat{\mathbf H}$ of the nuclear magnetic moment
$\hat{\boldsymbol\mu}=\mu\hat{\mathbf i}/i$ in the magnetic field produced by
the electron at the origin. The familiar formula of electrodynamics gives
this field as
\begin{equation}
 \hat{\mathbf H}=\frac1c\int\frac{\mathbf n\times\hat{\mathbf j}}{r^2}\,dV,
 \tag{121.8}
\end{equation}
where $\hat{\mathbf j}$ is the current-density operator produced by the
moving electron spin, while $\mathbf r=\mathbf n r$ is the radius vector from
the centre to the volume element $dV$\footnote{See II, \S~43, formula (43.7).
There the vector $\mathbf R$ points in the opposite direction, from $dV$ to
the centre (the point at which the field is observed).}. According to (115.4),
\[
 \hat{\mathbf j}=-2\mu_Bc\,\Curl{\psi^2\hat{\mathbf s}}
 =-2\mu_Bc\frac{d\psi^2(r)}{dr}\,\mathbf n\times\hat{\mathbf s}
\]
($\mu_B$ is the Bohr magneton). Writing $dV=r^2\,dr\,d\Omega$ and performing
the integrations, we find
\[
 \hat{\mathbf H}=-2\mu_B\int_0^\infty\frac{d\psi^2}{dr}\,dr
 \int\mathbf n\times(\mathbf n\times\hat{\mathbf s})\,d\Omega
 =-2\mu_B\psi^2(0)\frac{8\pi}{3}\hat{\mathbf s}.
\]
Thus the interaction operator is
\begin{equation}
 \hat V_{is}=-\hat{\boldsymbol\mu}\hat{\mathbf H}
 =\frac{16\pi}{3i}\mu\mu_B\psi^2(0)
 \hat{\mathbf i}\hat{\mathbf s}.
 \tag{121.9}
\end{equation}

If the total atomic angular momentum is $J=S=1/2$, the hyperfine splitting
produces a doublet, with $F=i\pm1/2$. Equations (121.5) and (121.9) give the
separation of its two levels as
\begin{equation}
 E_{i+1/2}-E_{i-1/2}
 =\frac{8\pi}{3i}\mu\mu_B(2i+1)\psi^2(0).
 \tag{121.10}
\end{equation}

\SourcePage{602}{605}
Since $\psi(0)$ is proportional to $\sqrt Z$ (see \S~71), this splitting
increases in proportion to the atomic number.

\ProblemHeading 1. Calculate the hyperfine splitting due to the magnetic interaction
for an atom having, outside closed shells, one electron with orbital angular
momentum $l$ (E.~Fermi, 1930).

\SolutionHeading Vector potential and magnetic-field intensity produced by
the nuclear magnetic moment $\boldsymbol\mu$ are
\[
 \mathbf A=\frac{\boldsymbol\mu\times\mathbf n}{r^2},\qquad
 \mathbf H=\frac{3\mathbf n(\boldsymbol\mu\mathbf n)-\boldsymbol\mu}{r^3}
\]
($\Div{\mathbf A}=0$). These expressions give the interaction operator in the
form
\[
 \frac{|e|}{mc}\mathbf A\hat{\mathbf p}
 +\frac{|e|\hbar}{mc}\hat{\mathbf H}\hat{\mathbf s}
 =\frac{2\mu_B}{r^3}\hat{\boldsymbol\mu}
 \left[\hat{\mathbf l}+3(\hat{\mathbf s}\mathbf n)\mathbf n
 -\hat{\mathbf s}\right].
\]
After averaging over a state of specified $j$, the expression in square
brackets is directed along $\mathbf j$. We may therefore write
\[
 \hat V_{ij}=2\mu_B(\hat{\boldsymbol\mu}\hat{\mathbf j})
 [\hat{\mathbf l}\hat{\mathbf j}
 +3(\hat{\mathbf s}\mathbf n)(\mathbf n\hat{\mathbf j})
 -\hat{\mathbf s}\hat{\mathbf j}]
 \frac{\overline{r^{-3}}}{j(j+1)}.
\]
The mean $\overline{n_in_k}$ was evaluated in the problem to \S~29. Using
that result and passing to eigenvalues, we obtain
\[
 \frac{2\mu_B\mu}{i}(\hat{\mathbf i}\hat{\mathbf j})
 \left[\hat{\mathbf l}\hat{\mathbf j}
 +\frac{2l(l+1)\hat{\mathbf s}\hat{\mathbf j}
 -6(\hat{\mathbf s}\hat{\mathbf l})(\hat{\mathbf j}\hat{\mathbf l})}
 {(2l-1)(2l+3)}\right]
 \frac{\overline{r^{-3}}}{j(j+1)},
\]
and a straightforward calculation finally gives
\[
 \frac{\mu_B\mu}{i}\frac{l(l+1)}{j(j+1)}F(F+1)\overline{r^{-3}},
\]
where $\mathbf F=\mathbf j+\mathbf i$ and $j=l\pm1/2$. The average of
$r^{-3}$ is taken over the radial part of the electron wavefunction.

\ProblemHeading 2. Determine the Zeeman splitting of the hyperfine-structure
components of an atomic level (S.~A.~Goudsmit, R.~F.~Bacher, 1930).

\SolutionHeading In formula (113.4), where the field is assumed so weak that
the splitting which it produces is small compared with the hyperfine
intervals, the averaging must now be carried out both over the electronic
state and over the orientations of the nuclear spin. The first averaging
gives $\Delta E=\mu_Bg_J\overline{J_z}H$, with the same $g_J$ as in (113.7).
By analogy with (113.5), the second averaging gives
\[
 \overline{J_z}=\frac{(\mathbf J\mathbf F)}{F^2}M_F.
\]
Consequently,

\SourcePage{603}{606}
\[
 \Delta E=\mu_Bg_FHM_F,\qquad
 g_F=g_J\frac{F(F+1)+J(J+1)-i(i+1)}{2F(F+1)}.
\]
